%! Diagonalizing a Matrix — Unit 6, Lesson 6.2 — Homework
\documentclass[10pt]{article}
\usepackage{linalg-article}
\usepackage{linalg-boxes}
\newcommand{\TallMath}[1]{$\displaystyle #1\rule[-1.4em]{0pt}{3.2em}$}
\newcommand{\xx}{\mathbf{x}}
\newcommand{\zero}{\mathbf{0}}

\begin{document}

\pageheader{Unit 6, Lesson 6.2}{Homework}
\namedateperiod

\vspace{0.12in}

\begin{notesbox}{Practice}
Line the eigenvectors into the columns of $X$ and the eigenvalues down the diagonal of $\Lambda$. Then
$A=X\Lambda X^{-1}$, and $A^k=X\Lambda^k X^{-1}$.

\begin{enumerate}[leftmargin=1.9em, itemsep=12pt]
  \item \textbf{Diagonalize.} For $A=\begin{bmatrix} 1 & 2 \\ 2 & 1 \end{bmatrix}$ (eigenvalues $\lambda=3,-1$
        with eigenvectors $\begin{bsmallmatrix}1\\1\end{bsmallmatrix},\begin{bsmallmatrix}1\\-1\end{bsmallmatrix}$),
        write $X$, $\Lambda$, and $X^{-1}$, then write $A=X\Lambda X^{-1}$. \writelines{3}
  \item \textbf{Powers made easy.} For $A=\begin{bmatrix} 2 & 3 \\ 0 & 1 \end{bmatrix}$ (eigenvalues
        $\lambda=2,1$ with eigenvectors $\begin{bsmallmatrix}1\\0\end{bsmallmatrix},\begin{bsmallmatrix}3\\-1\end{bsmallmatrix}$),
        use $A^4=X\Lambda^4 X^{-1}$ to compute $A^4$ \emph{without} multiplying $A$ four times. \writelines{3}
  \item \textbf{Diagonalizable or not?} For each matrix, decide whether it can be diagonalized and say why.
    \begin{enumerate}[label=(\alph*), leftmargin=1.8em, itemsep=4pt]
      \item $A=\begin{bmatrix} 3 & 0 \\ 1 & 3 \end{bmatrix}$ \hfill
      \item $A=\begin{bmatrix} 2 & 1 \\ 1 & 2 \end{bmatrix}$
    \end{enumerate}
    \vspace{1.2cm}
  \item \textbf{Reading powers off $\Lambda$.} A matrix has $A=X\Lambda X^{-1}$ with
        $\Lambda=\begin{bmatrix} 2 & 0 \\ 0 & 5 \end{bmatrix}$ (the eigenvectors don't change). Without
        computing $A$, state the eigenvalues of $A$, of $A^{3}$, and of $A^{-1}$. \writelines{2}
  \item \textbf{Explain.} Why does the factorization collapse to $A^k=X\Lambda^k X^{-1}$, and why must $A$
        have \emph{independent} eigenvectors for this to work? \writelines{2}
\end{enumerate}
\end{notesbox}

\begin{extensionbox}
\textbf{Back to the notes matrix.} In the notes you found $A^2$ for $A=\begin{bmatrix} 3 & 1 \\ 0 & 2 \end{bmatrix}$
(with $X=\begin{bsmallmatrix}1&1\\0&-1\end{bsmallmatrix}$, $X^{-1}=\begin{bsmallmatrix}1&1\\0&-1\end{bsmallmatrix}$).
Now use $A^4=X\Lambda^4 X^{-1}$ to compute $A^4$. \writelines{2}
\end{extensionbox}

\begin{spiralbox}
\textbf{Coming next --- Lesson 6.3: Symmetric Positive Definite Matrices.} You saw today that a matrix
diagonalizes only when it has enough independent eigenvectors. Next you'll meet a family that
\emph{always} does --- \textbf{symmetric} matrices --- and whose eigenvectors are even \emph{perpendicular},
giving the cleanest diagonalization of all.
\end{spiralbox}

\end{document}
